\chapter{Introduction to Agentic AI in Information Retrieval}
\label{ch:introduction}

\section{From Reactive Generation to Agency}

A generative language model is a function from a prompt to a completion. It
emits text and stops. It holds no goal that outlives the call, it cannot observe
what its output caused, and it has no way of concluding that what it produced
was not good enough. That describes a component, and it also describes most
systems currently called retrieval-augmented: embed a query, fetch a ranked
list, paste the list into a prompt, emit a completion. Retrieve once, generate
once. Nothing in that structure allows the behaviour to be anything but
reactive.

Agentic artificial intelligence names the class of systems built to remove that
limitation. An agentic system pursues a goal across several steps, chooses its
own actions, calls external instruments to carry them out, observes what came
back, and revises its behaviour in light of the
observation~\cite{zhang2024aiagent,wu2025agentic}. Four properties carry the
weight. \textbf{Autonomy}: the system, not the user, decides what happens next.
\textbf{Goal-directed planning}: the work is represented explicitly before it
is done, so the plan can be inspected, validated, repaired and discarded.
\textbf{Tool integration}: the model is not the source of facts but the
selector of instruments. \textbf{Adaptive reasoning}: the plan can be revised
after execution has begun, because an observation can invalidate it.

The fourth property is the difficult one and the one most often missing. A
system can be autonomous, planned and tool-using and still be a fixed pipeline
whose stages happen to be model calls. The industry touchstones the assignment
brief names, AutoGPT and Microsoft's Jarvis, are notable for attempting the
fourth property. The research literature shows the same direction. PaperQA
retrieves across scientific articles, judges what it found and synthesises from
the survivors~\cite{lala2023paperqa}. Agentic Reasoning adds tool-using
sub-agents and a mind-map agent~\cite{wu2025agentic}. CollEX explores
multimodal scientific collections~\cite{schneider2025collex}, agentic
contextual retrieval in telecommunications combines multi-source retrieval
with self-reflective validation~\cite{zhang2025toward}, and multi-agent
geospatial copilots decompose remote-sensing workflows across specialised
agents~\cite{lee2025multiagent}. What they share is the shape, and what they
leave open is the question this report exists to test: whether the feedback
edge, once actually implemented, earns the compute it consumes.

\section{The Limits of Static Ranking}

Classical information retrieval is the study of a scoring function. Given a
query and a corpus, rank the documents and return the top $k$. BM25 formalises
this for lexical matching~\cite{robertson2009probabilistic}, dense bi-encoders
replace term overlap with proximity in a learned space~\cite{xiao2023cpack},
reciprocal rank fusion combines rankings whose scores share no
scale~\cite{cormack2009reciprocal}, and BEIR measures how well any of it
transfers~\cite{thakur2021beir}. These are real and cumulative advances, and
every one of them improves the ordering of a single list produced from a
single query.

Consider a HotpotQA question~\cite{yang2018hotpotqa} that this report reuses
throughout:

\begin{quote}
\textit{When did the tour begin for the album Taylor Swift released on
October 22, 2012?}
\end{quote}

\noindent The passage \texttt{Red (Taylor Swift album)} says the album was
released on that date. The passage \texttt{The Red Tour} says the tour began on
March 13, 2013. Of the 66{,}581 passages in the processed HotpotQA corpus, none
contains both dates. That is not a statement about ranking quality. There is no
document for a ranker to place first, because no single document answers the
question. The bridge entity, the album title \textit{Red}, does not occur in
the question at all; it has to be derived before it can be searched for. A
lexical retriever cannot match a term that was never typed, and a dense
retriever embeds the whole question, so the sub-need ``identify the album'' is
diluted by the sub-need ``find its tour date'' rather than issued on its own.
This is the compositionality gap that Press et al.\ measure
directly~\cite{press2023selfask}, and it is what motivates interleaving
retrieval with reasoning rather than separating them~\cite{trivedi2023interleaving}.

Scale turns the difficulty into an impossibility. The two corpora hold
66{,}581 and 54{,}957 passages, so the number of ordered passage pairs is on
the order of $10^9$. The composition cannot be materialised in advance; it has
to be inferred from the question, which is a reasoning operation rather than a
ranking one. Two pressures compound it: the distractor setting fills each
candidate pool with passages chosen to look like the gold ones, and relevance
for the first hop is defined only relative to an intermediate goal the user
never stated.

\section{Problem Statement}

The problem is the design of a retrieval agent that, given a question and a
static corpus, does four things a ranker does not. It interprets an intent
whose surface form under-specifies the work. It decomposes that intent into
ordered sub-goals with explicit dependencies, so that one retrieval's result
can be substituted into the next query. It orchestrates the resulting workflow,
choosing an instrument for each sub-goal and traversing structure where the
corpus admits one. And it judges its own answer against the evidence it
actually retrieved before returning anything.

The fourth requirement is where a pipeline and an agent part company. The
system described here is, with one exception, a forward pipeline: plan,
execute, aggregate, synthesise. The exception is a single backward edge. A
Verifier scores each candidate answer against the evidence cited for it, and
when its confidence falls below a threshold it does not answer. It sends a
\texttt{ReplanDirective} back to the Planner carrying what could not be
established, which sub-queries returned nothing useful, and the text of every
sub-query already tried, which the next plan may not repeat. One edge is a
modest architectural change and a categorical behavioural one.

The system runs entirely on local hardware: \texttt{qwen3:8b} at Q4\_K\_M
quantisation through Ollama on a single RTX 5060 with 8~GB of VRAM, measured at
52 tokens per second warm, with no API key anywhere. At that throughput a
schema-constrained model call costs several seconds, so every call has to be
justified against a deterministic alternative. Tool selection is a decision
table over lexical features, entity linking is alias matching, graph traversal
is breadth-first search, and the confidence blend is arithmetic. A system that
cannot afford to think everywhere is forced to say where thinking pays, and, as
Chapter~\ref{ch:implementation} reports, the answer it arrived at was not the
one it was designed around.

\section{Objectives and Contributions}

The objective is a working agentic retrieval system covering all five tasks of
Table~1 of the assignment brief, whose central claim is tested rather than
asserted. The contributions are four.

\textbf{A four-agent architecture with full task coverage.} A Planner emits a
validated directed acyclic graph of sub-queries linked by placeholders. A
Retrieval Agent selects among BM25, dense and hybrid search under a conditional
reranking gate. A Knowledge Graph Navigator traverses an entity graph built from
corpus text alone. A Verifier performs self-reflective validation. An
orchestrator specified as an explicit state machine coordinates the four.

\textbf{The Verifier-to-Planner feedback loop, isolated by ablation.} A
configuration with the backward edge disabled and everything else held fixed
measures what the loop is worth, and the trace records which cycle produced the
selected answer.

\textbf{An evaluation protocol that prices agency.} Answer and retrieval quality
are reported beside model calls, tool calls, plan depth, re-plan rate and
latency, so no gain appears without its cost. Nine configurations run over two
multi-hop benchmarks, all eighteen cells on the same frozen 250 questions.

\textbf{A local, reproducible implementation, and an honest account of how
reproducible.} Every model runs on-device, prompts are versioned files whose
rendered hash is logged, and a test asserts that no module reaches the network.
Section~\ref{sec:reproducibility} reports the measured limit of that
reproducibility, which turned out to be a finding in its own right.

Table~\ref{tab:task-mapping} records where each task of the brief is
implemented and measured. One clause of the brief is deliberately not met: it
mentions orchestrating external APIs for real-time enrichment, and this system
exposes no external API at all, because a tool whose output changes between
runs makes the evaluation unrepeatable and the trace unfalsifiable. The
registry is written so that such a connector could be added as one more entry.

\begin{table}[htbp]
  \centering
  \small
  \caption{The five tasks of the brief's Table~1, the component implementing
  each, and where the report specifies and measures it. Four agents cover five
  tasks because query refinement and retrieval planning are one act performed
  twice: before any evidence exists, and again after the Verifier has said what
  is missing.}
  \label{tab:task-mapping}
  \begin{tabular}{p{4.0cm}p{4.6cm}p{5.6cm}}
    \toprule
    Task (brief, Table~1) & Agent or component & Specified in; measured in \\
    \midrule
    Autonomous query refinement &
      Planner: per-node rewrites; re-plan on a directive &
      \S\ref{sec:decomposition}, \S\ref{sec:verification};
      re-plan rate in Table~\ref{tab:agent-metrics-hotpotqa} \\
    Multi-step retrieval planning &
      Planner: validated sub-query DAG; orchestrator state machine &
      \S\ref{sec:decomposition}, \S\ref{sec:architecture};
      \S\ref{sec:ablations} \\
    Tool-augmented retrieval &
      Retrieval Agent: registry, routing table, reranking gate &
      \S\ref{sec:tools}; Table~\ref{tab:agent-metrics-hotpotqa} \\
    Knowledge graph traversal &
      KG Navigator: alias linking, bidirectional BFS &
      \S\ref{sec:kg-traversal}, \S\ref{sec:preprocessing};
      \S\ref{sec:ablations} \\
    Self-reflective validation and fact checking &
      Verifier: NLI entailment, citation grounding, feedback edge &
      \S\ref{sec:verification}; \S\ref{sec:ablations}, \S\ref{sec:errors} \\
    \bottomrule
  \end{tabular}
\end{table}

A fifth outcome was not designed but is reported because it is useful. During
development the 8B model produced correctly shaped decompositions in 8 of 8
cases with no format retries, and mislabelled its own plan's
\texttt{strategy} field in 7 of those 8. The system therefore derives the
strategy from the shape of the validated graph and discards what the model
claimed. Where a small model's self-report can be trusted, and where it cannot,
is a practical result for anyone building agents at this scale.

Chapter~\ref{ch:data} covers data and preprocessing, Chapter~\ref{ch:methodology}
the retrieval methodologies and why single-shot retrieval fails on compound
needs, Chapter~\ref{ch:implementation} the implementation, evaluation,
ablations, error analysis and the improvement loop that followed, and
Chapter~\ref{ch:conclusions} what the evidence supports and what follows.
